\documentclass[11pt,a4paper]{article}

\usepackage[margin=24mm,headheight=14pt]{geometry}
\usepackage[T1]{fontenc}
\usepackage[utf8]{inputenc}
\usepackage{lmodern}
\usepackage{xcolor}
\usepackage{tabularx}
\usepackage{booktabs}
\usepackage{enumitem}
\usepackage{fancyhdr}
\usepackage{hyperref}
\usepackage{needspace}

\definecolor{brand}{HTML}{3730A3}
\definecolor{ink}{HTML}{17243D}
\definecolor{muted}{HTML}{52627A}
\definecolor{pale}{HTML}{F1F3FF}

\hypersetup{
  colorlinks=true,
  linkcolor=brand,
  urlcolor=brand,
  pdftitle={Revenue Analytics Dashboard - Implementation Update},
  pdfauthor={Group 06}
}

\setlength{\parindent}{0pt}
\setlength{\parskip}{5pt}
\setlist[itemize]{leftmargin=18pt,itemsep=2pt,topsep=3pt}
\setlist[enumerate]{leftmargin=20pt,itemsep=2pt,topsep=3pt}
\renewcommand{\arraystretch}{1.2}

\pagestyle{fancy}
\fancyhf{}
\fancyhead[L]{\small\color{muted}Revenue Analytics Dashboard}
\fancyhead[R]{\small\color{muted}Group 06 -- Implementation update}
\fancyfoot[L]{\small\color{muted}Supersedes earlier status where inconsistent}
\fancyfoot[R]{\small\color{muted}Update \thepage}
\renewcommand{\headrulewidth}{0.3pt}
\renewcommand{\footrulewidth}{0.3pt}

\newcommand{\sectiontitle}[1]{%
  \Needspace{5\baselineskip}\vspace{2pt}\section*{\color{brand}#1}\vspace{-5pt}\hrule height 0.6pt\vspace{7pt}
}

\begin{document}

\begin{center}
  {\Large\bfseries\color{brand}REVENUE ANALYTICS DASHBOARD}\\[4pt]
  {\large Implementation and Verification Update}\\[7pt]
  {\color{muted}Group 06 \quad|\quad Web Design and Programming}\\
  {\color{muted}Status date: 10 November 2026}
\end{center}

\vspace{5pt}
\noindent\colorbox{pale}{\parbox{\dimexpr\linewidth-2\fboxsep\relax}{%
\textbf{Purpose of this addendum.} This update is appended to the report
\texttt{Group6\_Report\_LaTeX\_\_1\_.pdf}. The original PDF is retained unchanged
because its editable LaTeX source was not present in the project. This addendum
records the implemented state and supersedes earlier report statements wherever
they conflict with the current code, especially the database schema, chart
interaction, feedback, account functions, migration status, and test count.
}}

\sectiontitle{1. Current product scope}

Revenue Analytics is a browser-based mini BI application for students, small
sellers or teams, and analysts who have product or sales data in CSV files. It
supports dataset review and preparation, interactive charts, saved dashboards,
and forecasts evaluated against held-out historical data. It is not a live
market feed or an automated decision system.

\begin{itemize}
  \item \textbf{Data workspace:} validates CSV uploads (maximum 10 MiB), previews
  original and normalized cell values, recommends column types and roles, and
  reports missing values and exact duplicate rows. Common date, number,
  currency, and percentage formats are normalized where unambiguous. Removing
  duplicate or incomplete rows is an explicit user action; the importer does
  not silently delete them.
  \item \textbf{Visualize workspace:} supports Bar, Line, Pie, and Scatter charts,
  SUM, AVERAGE, COUNT, MIN, and MAX aggregation, dashboards with saved/reordered/
  removable visuals, and touch-friendly field dropdowns as an alternative to
  drag-and-drop.
  \item \textbf{Forecast workspace:} evaluates one selected numeric time series,
  optionally filtered to a product or group. It shows model comparison,
  held-out performance, future point estimates, data warnings, and saved runs.
\end{itemize}

\sectiontitle{2. User interface and interaction updates}

The sign-in page is also the bilingual project landing page. It now introduces
the project, its intended users, supported forecasting methods, and the purpose
and limitations of the application. The hero illustration is positioned below
the bold heading without rotation and scales to narrow screens. The landing
page includes a public feedback form.

On phones and tablets, chart fields can be assigned from dropdowns within each
field well; dragging is optional. When the dataset has usable ISO date values,
all chart types provide a searchable date checklist with \textit{Select all},
\textit{Select none}, and individual date selection. A selection immediately
rebuilds the active chart and is retained when the chart type changes. Saved
visuals persist the selected date column and dates, so the dashboard recreates
the same filtered view.

A Pie chart has one numeric measure and divides its aggregate into category
slices. This is appropriate for showing composition; Bar or Line charts are
available when comparing measures or trends. The interface explains this
distinction rather than suggesting that a pie chart should accept multiple
numeric measures.

The field-search input disables browser autocomplete and is cleared when the
Visualize page initializes. This prevents a stored email or other browser
suggestion from filtering the field list and making the dataset appear absent.
The selected CSV filename is shown before upload. Responsive layout was checked
in browser viewports from 320 px to 1440 px; this is viewport emulation, not a
claim of testing on every physical handset. Updated asset versions and
revalidation of served HTML reduce the chance of showing a cached older UI on
the first page load. Incorrect-password feedback for account deletion appears
beside the password field and clears when the user retries.

\sectiontitle{3. Accounts, privacy, and feedback}

There is no separate username: registration and sign-in use email. For the
OAuth2 password-form endpoint, the email is submitted in the field named
\texttt{username}. The backend trims and lowercases email addresses during
registration and sign-in, so email capitalization does not change account
identity. Passwords remain case-sensitive and are stored as bcrypt hashes.
Login tokens are JWTs signed using the required \texttt{SECRET\_KEY}.

Signed-in users can edit an optional display name and phone number, upload a
PNG/JPEG/WebP avatar up to 140 KiB, choose a per-browser light/dark/system
theme, and delete an account after password verification and confirmation.
Dataset ownership and roles are determined by the server.

Feedback is stored independently of accounts. Administrators can review it
through the authenticated \texttt{GET /api/v1/admin/feedback} endpoint; the
web interface does not currently include an inbox. A regular user cannot use
the administrator-only endpoint.

\sectiontitle{4. Data storage and current database schema}

The default database is SQLite at \texttt{data/app.db}; deployments can select
another database such as PostgreSQL through \texttt{DATABASE\_URL}. The server
stores parsed dataset rows, not a byte-for-byte copy of the uploaded CSV file.
Each dataset keeps original cell text, normalized analysis values, profiling
metadata, and the user's final semantic review. The browser keeps an active
working copy and token in \texttt{sessionStorage}; the database remains the
durable source of truth.

\begin{tabularx}{\linewidth}{@{}lX@{}}
\toprule
\textbf{Table} & \textbf{Stored information}\\
\midrule
\texttt{user} & Email, password hash, role, optional profile fields, avatar, and creation time.\\
\texttt{dataset} & Owner, filename, row and column information, original and normalized values, profiles, and semantic review.\\
\texttt{visual} & Dataset, chart type, field mapping, aggregation, optional \texttt{date\_filter} JSON, display order, and creation time. It does not store a chart image.\\
\texttt{forecastrun} & Owner and dataset, forecast settings, selected model/features, and saved result.\\
\texttt{feedback} & Submitter name, email, message, and creation time; readable only by administrators.\\
\bottomrule
\end{tabularx}

A visual is reconstructed from its dataset and saved chart settings. Deleting a
dataset cascades to its visuals and forecast runs. The optional date filter
records a date-column name and a unique list of selected \texttt{YYYY-MM-DD}
values; older visuals without a filter continue to display all dates.

\sectiontitle{5. Forecasting and interpretation}

Available methods are last-value Naive, the mean of up to the last seven
observations, damped Holt trend, optional Seasonal Naive when a season length is
specified, and optional Ridge regression when the user reviews and confirms
eligible historical numeric features. Ridge uses lagged target values and
selected historical feature values; it does not use actual values from inside
the forecast horizon. Imputation and scaling are fitted on training samples.

The forecast horizon is 1--90 periods. Two expanding chronological validation
windows select the method by pooled MAE. The final held-out test window is kept
separate from model selection; MAE and RMSE are reported, then the selected
method is refitted on usable history to produce future point estimates. Each
assisted run is saved for later review. The system does not currently provide
prediction intervals, guarantee future results, infer seasonality, or support
simultaneous multi-series forecasting.

\sectiontitle{6. Migrations, setup, and verification}

The current schema includes incremental Alembic revisions
\texttt{20261109\_05} (feedback storage) and \texttt{20261110\_06}
(optional saved-visual date filter). Revision \texttt{20261110\_06} depends on
\texttt{20261109\_05}; therefore, the complete
\texttt{migrations/versions/} directory must accompany the source code. For a
new or existing database, run \texttt{python -m alembic upgrade head} from the
project root before starting the application. Existing databases should be
backed up before schema migration. The application requires \texttt{SECRET\_KEY}
at startup.

On 10 November 2026, the backend suite completed with \textbf{135 passed}.
JavaScript syntax checks, \texttt{git diff --check}, browser checks for field
search and date filtering, and an isolated SQLite migration/startup check also
passed. The isolated checks used temporary databases; they did not replace
project data.

\sectiontitle{7. Limitations and follow-up}

CSV normalization uses heuristics, so ambiguous dates and locale-dependent
numbers require review. Missing periods are not automatically imputed for
forecasting. Forecasts are point estimates from historical patterns, not
business guarantees. Feedback review is API-only, and chart rendering depends
on the Chart.js CDN unless the library is bundled locally. Physical-device
accessibility, production deployment, and real-user forecast quality still
require separate evaluation.

\end{document}
